\documentclass{situationist-preprint}

\title{Narrative and Numerical Identity}
\msid{PID-09-b}
\msfamily{PID-09}
\msdiscipline{Personal Identity and the Self}
\mstype{Companion manuscript}
\mskeywords{narrative identity; numerical identity; characterization question; reidentification; dementia; co-authorship}

\begin{document}
\maketitle

\begin{abstract}
PID-09-a argues that narrative identity is real, co-authored and distinct from numerical identity. We develop the third claim. We argue that narrative identity presupposes numerical identity and cannot ground it; that the two vary independently, as cases of conversion, of dementia and of fission show; and that the confusion of the two explains both the appeal and the excesses of narrative theories of personal identity. The co-authorship of narrative identity, which PID-09-a establishes, is what allows it to survive the subject's loss of her own story while numerical identity continues unaffected.
\end{abstract}

\section{Two questions}

\citet{Schechtman1996} distinguished the \emph{reidentification question}, what makes a person at one time the same person as a person at another, from the \emph{characterization question}, which actions, experiences and traits are a person's own and what they mean in her life. PID-09-a, Definition~2.1, takes narrative identity as an answer to the second; Definition~2.3 takes numerical identity, with PID-01, as non-branching psychological continuity with bodily continuity as a tie-breaker, which answers the first. Lemma~3.3 of PID-09-a shows the two distinct by a difference in properties: narrative identity comes in degrees, is revisable and is not settled by its subject alone, while numerical identity is none of these.

\section{Presupposition}

\begin{proposition}\label{p:presup}
Narrative identity presupposes numerical identity: a narrative attribution of a past act to a person is correct only if the person is numerically identical with the agent of the act.
\end{proposition}

\begin{proof}
A narrative identity is a structure of attributions organized by an interpretation of a life as a temporally extended whole (PID-09-a, Definition~2.1). An attribution of an act to a person is correct only if she did it. A person did an act only if she is the one who performed it, which is a matter of numerical identity. A story that attributes to me the deeds of another is not a deep narrative identity of mine but a false story about me.
\end{proof}

Narrative theories that propose to answer the reidentification question by narrative therefore move in a circle. To say that a past action is mine because it belongs to my story requires that my story include only what I did, and what I did is fixed by a relation the story does not constitute.

\section{Independent variation}

The two relations also vary independently.

\paragraph{Conversion.} A person who undergoes religious conversion or recovery from addiction may revise the narrative of her life so thoroughly that past acts acquire new meanings, and some, once central, become peripheral. Her narrative identity changes greatly; her numerical identity does not change at all, since the person who converts is psychologically continuous with the one who did not.

\paragraph{Dementia.} In advanced dementia a person may lose the capacity to sustain a narrative of her own life. If narrative identity were authored by its subject alone, it would end. But it is co-authored (PID-09-a, Lemma~3.2): others who knew her continue to hold her story, to treat her in accordance with it and to keep in place the attributions that make her conduct intelligible \citep{Lindemann2014}. Her narrative identity is diminished but continues, sustained by others. Her numerical identity is a separate question, answered by the continuity of her remaining psychological capacities (PID-08) and, on PID-05, by the life of the organism she is. The two may come apart in either direction: others may hold a story she can no longer recognize, while numerical identity holds; and in the final stages, the story may outlast the psychological continuity it was about.

\paragraph{Fission.} In fission, two persons each continue the original's narrative. Before fission there was one story; afterwards each successor's story contains the same past. Narrative identity, which tolerates such sharing, does not tell us how many persons there are. Numerical identity does: there are two, and neither is the original (PID-03-a, Lemma~3.1).

\section{The appeal and the excess}

Narrative theories appeal because narrative identity answers the questions that matter most in practice: what we are responsible for, what our commitments mean, how others are to treat us \citep{MacIntyre1981}. Their excess lies in taking it to answer the metaphysical question as well, which yields the claim that those who do not live their lives as narratives lack a self, or are less fully persons. \citet{GStrawson2004} rightly rejects this claim for episodic people. The distinction defended here disarms it: an episodic person has a full numerical identity and a narrative identity largely authored by others, and neither fact makes her less a person.

\begin{objection}
The distinction concedes too much to the reidentification question. For practical purposes, characterization is all that matters, and numerical identity is a philosopher's abstraction.
\end{objection}

\begin{reply}
Characterization depends on reidentification by Proposition~\ref{p:presup}: responsibility, desert and the keeping of commitments all require that the one now held to account be the one who acted. What is true is that numerical identity carries little of what matters in survival (PID-06), while narrative identity carries much of what matters in the conduct of a life. That is a difference in weight, not in dependence.
\end{reply}

\section{Conclusion}

Narrative identity presupposes numerical identity, varies independently of it, and is held in part by others. The first fact rules out narrative accounts of reidentification; the second shows the two to be distinct; the third explains how a narrative identity can outlast its subject's capacity to sustain it.

\begin{thebibliography}{99}\small
\bibitem[Lindemann(2014)]{Lindemann2014} Lindemann, H. (2014). \emph{Holding and Letting Go: The Social Practice of Personal Identities}. Oxford University Press.
\bibitem[MacIntyre(1981)]{MacIntyre1981} MacIntyre, A. (1981). \emph{After Virtue: A Study in Moral Theory}. Duckworth.
\bibitem[Schechtman(1996)]{Schechtman1996} Schechtman, M. (1996). \emph{The Constitution of Selves}. Cornell University Press.
\bibitem[Strawson(2004)]{GStrawson2004} Strawson, G. (2004). Against narrativity. \emph{Ratio}, 17(4), 428--452.
\end{thebibliography}

\end{document}
